\chapter*{Introduction}
\addcontentsline{toc}{chapter}{Introduction}

L'objectif de ce mémoire est de présenter les travaux réalisés dans le cadre de mon stage de fin d'études, effectué au sein du laboratoire de recherche LATTICE\footnote{Le Lattice est un laboratoire du CNRS (UMR8094), rattaché à l’Ecole Normale Supérieure (IDEX PSL : Paris Sciences et Lettres) et à l’Université Sorbonne Nouvelle-Paris 3. Il regroupe des chercheurs CNRS, des enseignants-chercheurs de l’Université Paris Sorbonne Nouvelle et d’autres universités, ainsi que de nombreux autres personnels permanents ou non-permanents.
Les recherches du Lattice sont consacrées à la linguistique et au traitement automatique des langues.Dans le domaine linguistique proprement dit, nos recherches portent, d’une part, sur l’interface entre grammaire et lexique (dans le courant des grammaires cognitives), et d’autre part, sur le discours (en particulier la structuration du discours et la coréférence).
Le Lattice accorde également une place majeure à l’étude du changement linguistique, qu’il s’agisse de l’évolution du français – en particulier dans les domaines syntaxique et sémantique – ou des tendances plus générales à l’œuvre dans l’évolution des langues.
Dans le domaine du TAL, le Lattice explore les techniques neuronales depuis plusieurs années, et a récemment mis l’accent sur l’analyse syntaxique. Les travaux en linguistique et en TAL sont étroitement corrélés.
Enfin notre laboratoire est fortement impliqué dans le domaine des Humanités Numériques et contribue à la construction de ressources enrichies. https://www.lattice.cnrs.fr/ }. Au cours de ce stage, j’ai travaillé sur de nombreux aspects de \textit{TAL} (Traitement Automatique des Langues). Mon objectif principal était de déterminer comment était redigée l'ecriture naturaliste. Ce stage a été encadré par le professeur Dominque Legallois ainsi que par d'autre chercheurs du laboratoire à qui j'ai pu solicité leur aide notament lors d'utilisation de technique de \textit{machine Learning}.

La modélisation thématique, la stylométrie et la classification automatique constituent trois approches complémentaires de l’analyse quantitative des textes. Elles se distinguent principalement par leurs objectifs : la modélisation thématique cherche à faire émerger les structures sémantiques d’un corpus, la stylométrie mesure les caractéristiques formelles de l’écriture, tandis que la classification automatique attribue aux textes des catégories préalablement définies.

La modélisation thématique (\emph{topic modeling}) vise généralement à dégager, sans annotation préalable, les principaux thèmes présents dans un ensemble de documents. L’allocation de Dirichlet latente (\emph{Latent Dirichlet Allocation}, LDA) s’est imposée comme l’un des modèles probabilistes de référence. Elle représente chaque document comme un mélange de thèmes, eux-mêmes définis par des distributions de probabilités sur les mots du corpus.\footnote{David M. Blei, Andrew Y. Ng et Michael I. Jordan, « Latent Dirichlet Allocation », \emph{Journal of Machine Learning Research}, vol.~3, 2003, p.~993--1022, \href{https://www.jmlr.org/papers/v3/blei03a.html}{texte intégral en ligne}.} Les prolongements de cette méthode permettent notamment d’intégrer l’évolution temporelle des thèmes, leurs relations ou les métadonnées associées aux documents.\footnote{David M. Blei, « Probabilistic Topic Models », \emph{Communications of the ACM}, vol.~55, no~4, 2012, p.~77--84, \href{https://www.cs.columbia.edu/~blei/papers/Blei2012.pdf}{texte intégral en ligne}.} Plus récemment, les modèles neuronaux et les plongements contextuels ont renouvelé la modélisation thématique, en particulier pour l’analyse des textes courts, multilingues ou sémantiquement complexes.\footnote{Xiaobao Wu, Thong Nguyen et Anh Tuan Luu, « A Survey on Neural Topic Models: Methods, Applications, and Challenges », \emph{Artificial Intelligence Review}, vol.~57, article~18, 2024, \href{https://doi.org/10.1007/s10462-023-10661-7}{DOI et texte intégral}.} BERTopic, par exemple, combine des représentations produites par un modèle de langue, une méthode de regroupement automatique et une variante du TF--IDF pour caractériser les thèmes obtenus.\footnote{Maarten Grootendorst, « BERTopic: Neural Topic Modeling with a Class-Based TF-IDF Procedure », prépublication arXiv, 2022, \href{https://arxiv.org/abs/2203.05794}{texte intégral en ligne}.} L’interprétation des résultats demeure néanmoins délicate : les thèmes dépendent de la constitution et du prétraitement du corpus, du nombre de thèmes demandé et des paramètres du modèle. Par ailleurs, une bonne performance statistique ne garantit pas nécessairement que les thèmes produits soient cohérents pour un lecteur humain.\footnote{Jonathan Chang, Jordan Boyd-Graber, Sean Gerrish, Chong Wang et David M. Blei, « Reading Tea Leaves: How Humans Interpret Topic Models », dans \emph{Advances in Neural Information Processing Systems 22}, 2009, p.~288--296, \href{https://papers.neurips.cc/paper/2009/hash/f92586a25bb3145facd64ab20fd554ff-Abstract.html}{texte intégral en ligne}.}

La stylométrie cherche, pour sa part, à mesurer les régularités formelles d’une écriture à partir de caractéristiques quantifiables, telles que les fréquences lexicales, l’emploi des mots-outils, la richesse du vocabulaire, la longueur des phrases, la ponctuation, les catégories grammaticales ou les séquences de caractères. Elle a d’abord été principalement utilisée pour résoudre des problèmes d’attribution d’auteur, comme l’illustre l’étude fondatrice de Frederick Mosteller et David Wallace sur les articles anonymes des \emph{Federalist Papers}.\footnote{Frederick Mosteller et David L. Wallace, \emph{Inference and Disputed Authorship: The Federalist}, Reading, Addison-Wesley, 1964 ; réédition avec une introduction de John Nerbonne, Stanford, CSLI Publications, 2007, \href{https://web.stanford.edu/group/cslipublications/cslipublications/site/1575865521.shtml}{notice éditoriale}.} Parmi les méthodes classiques, le \emph{Delta} de John Burrows mesure la distance stylistique entre plusieurs textes à partir des fréquences standardisées de leurs mots les plus courants.\footnote{John Burrows, « `Delta': A Measure of Stylistic Difference and a Guide to Likely Authorship », \emph{Literary and Linguistic Computing}, vol.~17, no~3, 2002, p.~267--287, \href{https://doi.org/10.1093/llc/17.3.267}{DOI}.} Les recherches contemporaines mobilisent également les n-grammes de caractères, les méthodes d’apprentissage automatique et les représentations produites par les modèles de langue.\footnote{Efstathios Stamatatos, « A Survey of Modern Authorship Attribution Methods », \emph{Journal of the American Society for Information Science and Technology}, vol.~60, no~3, 2009, p.~538--556, \href{https://www.clips.uantwerpen.be/~walter/educational/material/Stamatatos_survey2009.pdf}{texte intégral en ligne}.} La stylométrie comprend désormais plusieurs tâches, notamment l’attribution, la vérification et le profilage d’auteur, ainsi que l’étude de l’évolution du style au cours du temps.\footnote{Enzo Terreau, \emph{Apprentissage de représentations d’auteurs et d’autrices à partir de modèles de langue pour l’analyse des dynamiques d’écriture}, thèse de doctorat en informatique, Université Lyon~2, 2024, \href{https://theses.fr/2024LYO20001}{notice et résumé en ligne}.} Ses résultats restent toutefois sensibles à la longueur des documents et aux effets du genre, du thème, de la période historique, du support ou de la traduction. Ces facteurs doivent donc être contrôlés afin de ne pas les confondre avec le style individuel.

Enfin, la classification automatique consiste à attribuer un texte à une ou plusieurs catégories définies à l’avance. Contrairement à la modélisation thématique, elle relève généralement de l’apprentissage supervisé : le modèle apprend à partir d’un corpus dont les documents ont préalablement été annotés.\footnote{Fabrizio Sebastiani, « Machine Learning in Automated Text Categorization », \emph{ACM Computing Surveys}, vol.~34, no~1, 2002, p.~1--47, \href{https://arxiv.org/abs/cs/0110053}{texte intégral en ligne}.} Les approches classiques représentent les documents par des fréquences lexicales ou des pondérations TF--IDF, puis utilisent des méthodes comme le classifieur bayésien naïf, la régression logistique ou les machines à vecteurs de support.\footnote{Christopher D. Manning, Prabhakar Raghavan et Hinrich Schütze, \emph{Introduction to Information Retrieval}, Cambridge, Cambridge University Press, 2008, chap.~13 à 15, \href{https://nlp.stanford.edu/IR-book/information-retrieval-book.html}{édition intégrale en ligne}.} Ces méthodes ont été complétées par les réseaux de neurones, puis par les modèles de langue préentraînés fondés sur l’architecture Transformer. BERT a notamment popularisé une procédure dans laquelle un modèle préentraîné sur de grandes quantités de textes est ensuite spécialisé sur une tâche de classification à partir d’un corpus annoté plus restreint.\footnote{Jacob Devlin, Ming-Wei Chang, Kenton Lee et Kristina Toutanova, « BERT: Pre-training of Deep Bidirectional Transformers for Language Understanding », dans \emph{Proceedings of NAACL-HLT 2019}, p.~4171--4186, \href{https://aclanthology.org/N19-1423/}{texte intégral en ligne}.} Les grands modèles de langue permettent également d’effectuer certaines classifications sans apprentissage spécifique ou à partir d’un nombre limité d’exemples. Ils ne rendent cependant pas systématiquement obsolètes les modèles spécialisés : pour de nombreuses tâches, un modèle plus petit correctement entraîné peut rester plus performant, plus stable et moins coûteux.\footnote{Aleksandra Edwards et Jose Camacho-Collados, « Language Models for Text Classification: Is In-Context Learning Enough? », dans \emph{Proceedings of LREC-COLING 2024}, p.~10058--10072, \href{https://aclanthology.org/2024.lrec-main.879/}{texte intégral en ligne}.}

Dans ces trois domaines, le choix de la méthode ne peut être séparé des conditions de constitution et d’annotation du corpus. Une évaluation rigoureuse suppose notamment de contrôler les effets d’auteur, de source, de genre et de période, de séparer strictement les données d’apprentissage et de test, et de confronter les résultats quantitatifs à une lecture qualitative des textes.\footnote{Justin Grimmer et Brandon M. Stewart, « Text as Data: The Promise and Pitfalls of Automatic Content Analysis Methods for Political Texts », \emph{Political Analysis}, vol.~21, no~3, 2013, p.~267--297, \href{https://doi.org/10.1093/pan/mps028}{DOI}.}
 
Mon obje
Le naturalisme constitue l’un des principaux mouvements littéraires et artistiques de la seconde moitié du XIXe siècle. Prolongeant certaines ambitions du réalisme, il se donne pour objectif de représenter l’être humain et la société avec une précision inspirée des sciences expérimentales. L’œuvre littéraire ne repose donc plus seulement sur l’imagination de l’écrivain : elle s’appuie sur l’observation du réel, sur une documentation minutieuse ainsi que sur l’étude des mécanismes physiologiques, sociaux et psychologiques qui déterminent les comportements individuels.

Émile Zola\footnote{Pour ce qui est des informations sur Émile Zola, je me suis grandement référé à sa page Wikipédia, qui est considérée comme une page de qualité : \href{{https://fr.wikipedia.org/wiki/\%C3\%89mile_Zola}}{lien de sa page.}}, principal théoricien du mouvement, définit ainsi le naturalisme comme « l’esprit scientifique porté dans toutes nos connaissances ». Il entend substituer à « l’homme métaphysique » un « homme physiologique », inséparable de son hérédité et du milieu dans lequel il évolue\footnote{\url{https://fr.wikipedia.org/wiki/Naturalisme_(litt\%C3\%A9rature)}}. Inspiré notamment par \textit{l’Introduction à l’étude de la médecine expérimentale} de Claude Bernard, Zola conçoit le romancier comme un observateur et un expérimentateur. Celui-ci commence par recueillir et consigner des faits, puis place ses personnages dans des circonstances déterminées afin d’étudier le fonctionnement de leurs passions, de leurs instincts et de leurs conduites. Le roman devient alors une forme d’expérience permettant de mettre au jour les lois qui régissent l’individu et la société\footnote{Émile Zola, \textit{Le Roman expérimental}, Paris, Charpentier, 1881.}.

Cette conception repose principalement sur le déterminisme. Dans la perspective naturaliste, la personnalité et la destinée d’un individu sont largement conditionnées par son héritage physiologique, son environnement familial, sa situation économique et son appartenance sociale. Le naturalisme ne doit cependant pas être réduit à un ensemble de procédés stylistiques\footnote{\url{https://fr.wikipedia.org/wiki/Naturalisme_(litt\%C3\%A9rature)}}. Paul Alexis le présente avant tout comme une méthode permettant de penser\footnote{Jules Huret, \textit{Enquête sur l'évolution littéraire}, Paris, Charpentier, 1891.}, de voir, d’étudier et d’expérimenter. De même, Camille Lemonnier l’envisage comme un « réalisme agrandi » par l’analyse des milieux et par l’apport des sciences naturelles et sociales. Le mouvement s’inscrit ainsi dans une vision matérialiste et empirique du monde, selon laquelle les phénomènes humains doivent être expliqués à partir de causes observables\footnote{\url{https://fr.wikipedia.org/wiki/Naturalisme_(litt\%C3\%A9rature)}}.

Sur le plan esthétique et narratif, cette ambition scientifique entraîne une remise en cause des conventions romanesques traditionnelles. Les naturalistes rejettent les intrigues artificielles, les péripéties invraisemblables et les personnages idéalisés. À la construction d’une fable spectaculaire, ils préfèrent la restitution d’une « tranche de vie », c’est-à-dire l’étude détaillée d’un fragment d’existence ordinaire. Le roman prend alors la forme d’une monographie consacrée à un individu, à une famille, à un métier ou à un milieu social. Le narrateur tend à s’effacer derrière les faits et adopte une posture d’impersonnalité comparable à celle d’un greffier chargé d’exposer les pièces d’un dossier sans formuler explicitement de jugement\footnote{Émile Zola, \textit{Le Roman expérimental}}.

La description occupe, dans cette perspective, une fonction essentielle. Elle ne constitue pas un simple ornement esthétique, mais un instrument d’analyse et d’explication. Les logements, les vêtements, les paysages, les lieux de travail et les conditions matérielles d’existence contribuent à façonner les personnages. Décrire un environnement revient donc à rendre visibles les forces qui agissent sur l’individu. Cette attention portée aux déterminations sociales conduit les écrivains naturalistes à s’intéresser aux catégories jusque-là marginalisées par la littérature : les ouvriers, les paysans, les prostituées, les criminels, les malades ou les exclus. Le héros traditionnel s’efface au profit de personnages ordinaires, souvent confrontés à la pauvreté, à la violence, à la maladie ou au déclassement.

La représentation de ces réalités a toutefois suscité d’importantes controverses morales et esthétiques. La crudité de certaines scènes, l’étude des pathologies et l’attention accordée aux formes de déchéance ont conduit les adversaires du mouvement à l’accuser de complaisance envers la laideur et le vice. Zola défend au contraire la dimension morale de sa démarche. À ses yeux, le romancier agit comme un anatomiste ou un chirurgien : en analysant les causes matérielles et sociales des comportements, il contribue à mieux comprendre les dysfonctionnements de la société et fournit les connaissances nécessaires à leur éventuelle correction. L’écrivain naturaliste serait ainsi un « moraliste expérimentateur », dont la fonction n’est pas de condamner, mais d’expliquer\footnote{\textit{Ibid.}}.

Cette conception est notamment contestée par Ferdinand Brunetière\footnote{Ferdinand Brunetière, \textit{Le roman naturaliste}, Paris, Calmann-Lévy, 1883.}, qui reproche au naturalisme français de privilégier la recherche de la grossièreté, le pessimisme et l’« écriture artiste » au détriment de la profondeur psychologique. Il oppose à ce qu’il considère comme un « faux naturalisme » une littérature incarnée par Dickens, George Eliot ou Tolstoï, caractérisée par la probité de l’observation, la simplicité de l’exécution et la sympathie envers les humbles. Edmond Cattier insiste également sur le rôle de l’« imagination sympathique »\footnote{Edmond Cattier, \textit{Le Naturalisme littéraire}, Bruxelles, Weissenbruch, 1897.}, faculté par laquelle l’écrivain dépasse la simple observation extérieure pour pénétrer la conscience d’autrui, comprendre les motifs de ses actions et restituer la complexité de son expérience. Ces critiques mettent en évidence une tension centrale du naturalisme : celle qui oppose l’ambition d’objectivité scientifique à la nécessité d’une compréhension sensible et psychologique de l’être humain.

Le naturalisme ne forme d’ailleurs pas une école parfaitement homogène. S’il revendique l’héritage de Diderot, de Rousseau, de Balzac, de Stendhal, de Flaubert et des frères Goncourt, ses représentants divergent sur la définition même du mouvement. Edmond de Goncourt préfère parler de « naturisme », Joris-Karl Huysmans d’« intimisme », Louis-Marie Desprez d’« impressionnisme », tandis que Guy de Maupassant propose le terme d’« illusionnistes ». Cette diversité terminologique révèle l’existence de pratiques et de conceptions distinctes derrière l’apparente unité du groupe naturaliste\footnote{\url{https://fr.wikipedia.org/wiki/Naturalisme_(litt\%C3\%A9rature)}}.

Le succès des \textit{Soirées de Médan}, publiées en 1880, marque l’un des moments les plus emblématiques de cette aventure collective. Toutefois, l’unité du mouvement se fragilise rapidement. Le \textit{Manifeste des Cinq}, publié en 1887 contre \textit{La Terre} de Zola, rend visibles les désaccords internes, tandis que plusieurs auteurs s’éloignent progressivement de l’esthétique naturaliste. Huysmans, notamment, ouvre avec \textit{À rebours} la voie à une littérature décadente et symboliste\footnote{\textit{Ibid.}}. À la fin du siècle, le développement du roman psychologique et du symbolisme contribue ainsi au déclin du naturalisme en tant qu’école constituée. La formule attribuée à Stéphane Mallarmé, selon laquelle « le mot mourra quand Zola aura achevé son œuvre », souligne à la fois le rôle déterminant de Zola dans la définition du mouvement et la difficulté du naturalisme à lui survivre sous une forme unifiée\footnote{Jules Huret, \textit{Enquête sur l'évolution littéraire}, Paris, Charpentier, 1891.}.

Dans quelle mesure les méthodes de modélisation thématique et de classification automatique permettent-elles d’identifier des marqueurs du naturalisme sans confondre signature d’auteur, contenu thématique et artefacts de constitution du corpus ?

Afin de répondre à cette problématique, ce mémoire s’organise en trois chapitres complémentaires. Le premier présente une exploration thématique classique des trente et un romans d’Émile Zola à l’aide de la factorisation en matrices non négatives et de la méthode K-means. Cette première analyse vise à dégager des régularités lexicales et thématiques internes à l’œuvre zolienne, considérées comme des marqueurs candidats plutôt que comme des preuves immédiates du naturalisme. Le deuxième chapitre étudie ces régularités au moyen de BERTopic et analyse leur évolution au cours de la carrière de Zola. Le troisième chapitre confronte les hypothèses issues de ces analyses à un corpus comparatif grâce à une classification automatique opposant romans naturalistes et romans relevant d’autres courants. Une attention particulière est accordée aux risques de confusion entre signature d’auteur, contenu thématique et artefacts liés à la constitution du corpus. Enfin, les choix de nettoyage, l’organisation de la base de données et les principaux éléments nécessaires à la reproductibilité des expériences sont présentés en annexe.





